\documentclass[12pt,a4paper]{article}
\usepackage[utf8]{inputenc}
\usepackage[spanish]{babel}
\usepackage{geometry}
\geometry{top=3cm, bottom=3cm, left=2.5cm, right=2.5cm}
\usepackage{graphicx}
\usepackage{booktabs}
\usepackage{xcolor}
\usepackage{listings}
\usepackage{hyperref}
\usepackage{fancyhdr}
\usepackage{amsmath}
\usepackage{float}

% --- Configuracion de Colores Estilo Code ---
\definecolor{codegreen}{rgb}{0,0.6,0}
\definecolor{codegray}{rgb}{0.5,0.5,0.5}
\definecolor{codeblue}{rgb}{0.1,0.4,0.8}
\definecolor{backcolour}{rgb}{0.96,0.96,0.96}

\lstdefinestyle{mystyle}{
    backgroundcolor=\color{backcolour},   
    commentstyle=\color{codegreen},
    keywordstyle=\color{codeblue},
    numberstyle=\tiny\color{codegray},
    stringstyle=\color{orange},
    basicstyle=\ttfamily\footnotesize,
    breakatwhitespace=false,         
    breaklines=true,                 
    captionpos=b,                    
    keepspaces=true,                 
    numbers=left,                    
    numbersep=5pt,                  
    showspaces=false,                
    showstringspaces=false,
    showtabs=false,                  
    tabsize=2
}
\lstset{style=mystyle}

% --- Configuracion de Encabezado y Pie ---
\pagestyle{fancy}
\fancyhf{}
\lhead{Equipo Linus - Logica Simbolica}
\rhead{Informe de Proyecto}
\cfoot{\thepage}

\begin{document}

% ==================== PORTADA ====================
\begin{titlepage}
    \centering
    \vspace*{2cm}
    {\scshape\LARGE Universidad Nacional Pedro Ruiz Gallo \par}
    \vspace{1cm}
    {\scshape\Large Facultad de Ingenieria Civil, de Sistemas y de Arquitectura \par}
    \vspace{0.5cm}
    {\large Escuela Profesional de Ingenieria de Sistemas \par}
    \vspace{1.5cm}
    
    {\Huge\bfseries Sistema Web Interactivo de Teoria de Conjuntos y Diagramas de Venn \par}
    \vspace{1cm}
    {\Large\bfseries Informe Tecnico Final del Proyecto \par}
    \vspace{2cm}
    
    {\large\bfseries Integrantes - Equipo Linus (Grupo 4):\par}
    \vspace{0.3cm}
    {\large
    Jordy Jose (Sublider) \\
    Junior \\
    Mike \\
    Sergio \\
    Fernando \\
    Alejandro \par}
    
    \vspace{2cm}
    {\large\bfseries Docente:\par}
    {\large Dr. Mardo Victor Gonzales Herrera \par}
    
    \vfill
    {\large Lambayeque, 2026 \par}
\end{titlepage}

\tableofcontents
\newpage

% ==================== SECCION 1 ====================
\section{Introduccion y Propuesta del Proyecto (Sprint 1)}
El aprendizaje de la Teoria de Conjuntos representa una base fundamental en la logica matematica y la ciencia de la computacion. No obstante, la visualizacion espacial de las operaciones (como uniones, intersecciones complejas y complementos) y el calculo manual del Conjunto Potencia $2^N$ suelen ser fuentes de confusion academica.

Para mitigar este problema, el \textbf{Equipo Linus} desarrollo una herramienta de software interactiva y didactica basada en la web. La funcionalidad se diseno a partir de bocetos y maquetas preliminares de la interfaz que sirvieron de base para estructurar el software final.

\begin{figure}[H]
    \centering
    % Descomentar y definir el archivo de imagen cuando lo tengas
    % \includegraphics[width=0.7\textwidth]{boceto_propuesta.png}
    \caption{Diseno en maqueta inicial (Wireframe) del proyecto.}
    \label{fig:boceto}
\end{figure}

% ==================== SECCION 2 ====================
\section{Arquitectura del Software y Flujo de Trabajo (Git Flow)}
El proyecto se diseño bajo una arquitectura desacoplada para asegurar escalabilidad y modularidad, separando la logica pura (el motor) de la interfaz de usuario. 

\subsection{Estructura de Directorios}
La organizacion de los archivos creados en el repositorio es la siguiente:
\begin{itemize}
    \item \textbf{Capa Logica:} \texttt{src/lib/sets/operations.ts}
    \item \textbf{Pruebas Unitarias:} \texttt{src/lib/sets/operations.test.ts}
    \item \textbf{Capa de Interfaz:} \texttt{src/pages/conjuntos/index.vue}
    \item \textbf{Documentacion de Avance:} \texttt{TAREAS/02\_Motor/linus/avance.md}, \texttt{TAREAS/03\_UI/linus/avance.md}, \texttt{TAREAS/04\_Deploy/linus/uso.md}
\end{itemize}

\subsection{Flujo de Git por Ramas}
Para mantener el repositorio limpio, cada fase se desarrollo en una rama de Git independiente antes de enviarse al repositorio principal mediante Pull Requests (PR):
\begin{itemize}
    \item \texttt{equipo-4-linus-propuesta} (Sprint 1)
    \item \texttt{equipo-4-linus-motor} (Sprint 2)
    \item \texttt{equipo-4-linus-ui} (Sprint 3)
    \item \texttt{equipo-4-linus-deploy} (Sprint 4)
\end{itemize}

\begin{figure}[H]
    \centering
    % \includegraphics[width=0.8\textwidth]{captura_ramas_git.png}
    \caption{Historial de desarrollo secuencial por ramas del Equipo Linus.}
    \label{fig:ramas}
\end{figure}

% ==================== SECCION 3 ====================
\section{Capa Logica: El Motor Matematico (Sprint 2)}
El motor logico realiza todos los calculos matematicos de conjuntos. Se implemento en TypeScript de forma estricta (tipado estricto) utilizando la estructura de datos nativa \texttt{Set} de ES6.

\subsection{Ventaja del Tipo Set}
En lugar de arreglos comunes, se seleccionaron \texttt{Set} debido a que:
\begin{enumerate}
    \item Filtran automaticamente elementos repetidos.
    \item La busqueda de elementos con \texttt{.has()} tiene una complejidad temporal de $O(1)$ (tiempo constante) gracias al uso de tablas hash internas.
\end{enumerate}

\subsection{Codigo Fuente del Motor Logico}
A continuacion, se detalla el codigo matematico de \texttt{src/lib/sets/operations.ts}:

\begin{lstlisting}[language=TypeScript, caption=Codigo de operaciones de conjuntos]
// src/lib/sets/operations.ts

export function union<T>(a: Set<T>, b: Set<T>): Set<T> {
  return new Set([...a, ...b]);
}

export function interseccion<T>(a: Set<T>, b: Set<T>): Set<T> {
  return new Set([...a].filter(x => b.has(x)));
}

export function diferencia<T>(a: Set<T>, b: Set<T>): Set<T> {
  return new Set([...a].filter(x => !b.has(x)));
}

export function complemento<T>(universo: Set<T>, a: Set<T>): Set<T> {
  return diferencia(universo, a);
}

export function potencia<T>(a: Set<T>): Set<Set<T>> {
  const elementos = Array.from(a);
  const subconjuntos = new Set<Set<T>>();
  const numSubconjuntos = Math.pow(2, elementos.length);
  for (let i = 0; i < numSubconjuntos; i++) {
    const subconjunto = new Set<T>();
    for (let j = 0; j < elementos.length; j++) {
      if ((i & (1 << j)) !== 0) {
        subconjunto.add(elementos[j]);
      }
    }
    subconjuntos.add(subconjunto);
  }
  return subconjuntos;
}
\end{lstlisting}

% ==================== SECCION 4 ====================
\section{Capa Visual: Interfaz con Vibe Duolingo (Sprint 3)}
La pantalla interactiva en \texttt{index.vue} conecta el motor logico con el usuario final de manera atractiva y sencilla, adoptando el diseno de estilo ``Duolingo'' (color azul de acento, bordes redondeados amplios y tarjetas limpias).

\subsection{Reactividad en Vue 3}
\begin{itemize}
    \item \textbf{\texttt{ref}:} Usado para registrar el texto ingresado por el usuario (Inputs del Universo, A y B).
    \item \textbf{\texttt{computed}:} Propiedades calculadas en memoria con cache de procesamiento. Permiten actualizar el resultado matematico y el diagrama SVG automaticamente en tiempo real solo cuando cambian los inputs de origen.
\end{itemize}

\subsection{Iluminacion Dinamica del Diagrama de Venn (SVG)}
En lugar de importar una libreria pesada de graficos, se programo un diagrama vectorial SVG nativo. El color de las regiones cambia en tiempo real mediante la clase reactiva \texttt{vennColores}, asignando opacidades al circulo de la izquierda, derecha o a la interseccion segun la operacion seleccionada.

\begin{figure}[H]
    \centering
    % \includegraphics[width=0.8\textwidth]{captura_interfaz.png}
    \caption{Captura de pantalla de la interfaz interactiva y el Diagrama de Venn en accion.}
    \label{fig:interfaz}
\end{figure}

% ==================== SECCION 5 ====================
\section{Pruebas de Calidad del Codigo (Sprint 4)}
Para asegurar que las matematicas de la aplicacion sean infalibles y no presenten errores logicos, se diseno una bateria de 14 pruebas unitarias con la suite de testing \texttt{Vitest}.

Estas pruebas comprueban leyes matematicas fundamentales, tales como:
\begin{itemize}
    \item Ley conmutativa: $A \cup B = B \cup A$
    \item Ley distributiva: $A \cap (B \cup C) = (A \cap B) \cup (A \cap C)$
\end{itemize}

\begin{figure}[H]
    \centering
    % \includegraphics[width=0.7\textwidth]{captura_pruebas.png}
    \caption{Ejecucion exitosa de la bateria de pruebas logicas.}
    \label{fig:tests}
\end{figure}

% ==================== SECCION 6 ====================
\section{Participacion y Aportes del Equipo}
La distribucion de responsabilidades en el desarrollo de los entregables se resume en la siguiente tabla:

\begin{table}[H]
    \centering
    \caption{Participacion y roles de cada integrante del equipo.}
    \begin{tabular}{lll}
        \toprule
        \textbf{Integrante} & \textbf{Rol Principal} & \textbf{Contribucion Clave} \\
        \midrule
        Jordy Jose & Sublider / Git Manager & Gestion de ramas, fusiones y desarrollo de UI. \\
        Junior & Analista de Logica & Verificacion del modulo de propiedades y pertenencia. \\
        Mike & Ingeniero de Pruebas & Implementacion de casos limite en Vitest. \\
        Sergio & Desarrollador Frontend & Implementacion de la visualizacion SVG del diagrama. \\
        Fernando & Documentador Tecnico & Elaboracion del manual de usuario y documentacion. \\
        Alejandro & Disenador de Software & Modelado inicial del comportamiento del sistema. \\
        \bottomrule
    \end{tabular}
\end{table}

% ==================== SECCION 7 ====================
\section{Conclusiones}
\begin{enumerate}
    \item Se logro construir una aplicacion web interactiva que traduce conceptos matematicos abstractos de la teoria de conjuntos en graficos dinamicos comprensibles.
    \item La separacion limpia entre la logica matematica (\texttt{operations.ts}) y la presentacion visual (\texttt{index.vue}) demostro ser clave para el desarrollo colaborativo y la ejecucion de pruebas robustas sin acoplamiento grafico.
\end{enumerate}

\end{document}
